Using the normalized standard resolution for $\Z$, we apply the functor $\Hom_{\Z G}(-,A)$ to obtain the cochain complex $\Hom(P_*/D_*, A)$ which will allow us to make the necessary computation. The relevant part of the complex is given below:
\[
A\xlra{d_0}\Hom(G,A)\xlra{d_1}\Hom(G^2,A)
\]
We know that $H^1(G;A)=\ker(d_1)/\im(d_0)$, so we simply calculate this quotient. We can see immediately that $d_0\equiv 0$, since by definition, $d_0(a)=ga-a$, but since $G$ acts on $A$ trivially, we know that $ga=a$ and thus $d_0(a)=0$. Now, a 1-cocycle is defined by the fact that $d_1f=0$, which means that $d_1f(g_1,g_2)= g_1 \cdot f(g_2) - f(g_1 g_2) + f(g_1)=0$, and so $f(g_1g_2)=f(g_1)+f(g_2)$, again by the triviality of the $G$ action. This means that $f$ is a homomorphism and we are done.

\begin{question}
    
\end{question}
We will first check that this action of $S$ is well-defined, i.e. that $s\cdot f\in\Hom_R(S,M)$. We check:
\[
s\cdot f(s_1+s_2)=f((s_1+s_2)\cdot s)=f(s_1\cdot s)+f(s_2\cdot s)=s\cdot f(s_1)+s\cdot f(s_2)
due to the fact that $f\in\Hom_R(S,M)$. 
Similarly,
\[
r\cdot (s\cdot f)(s')=f((r\cdot s')\cdot s)=f(\a(r)s's)=f((r\cdot s')\cdot s)=s\cdot f(r\cdot s')
\]
so that $f$ is $R$-linear. This proves that this action is well-defined. Now, we check that action of $S$ makes $\Hom_R(S,M)$ into an $S$-module. Let $f,g\in\Hom_R(S,M)$. We check:
\[
(s_1+s_2)\cdot f(s')=f((s_1+s_2)\cdot s')=f(s_1\cdot s')+f(s_2\cdot s')=s_1\cdot f(s')+s_2\cdot f(s')=(s_1+s_2)\cdot f(s')
\]
and
\[
s\cdot (f+g)(s')=s\cdot f(s')+s\cdot g(s')=(s\cdot f)(s')+(s\cdot g)(s')
\]
so that $\Hom_R(S,M)$ is an $S$-module.


\begin{question}

\end{question}
In what follows given an element $x\in \Z G$, we write $x=\sum_i a_i t^i$ as the general expansion of $x$. 

\begin{enumerate}
\item[(a)] If $x\in \Z G$ is in the kernel of the augmentation map $\e$, then we write $x=\sum_i a_i g_i$ where $\sum_i a_i =0$. We then notice that $x=\sum_i a_i g_i-\sum_i a_i = \sum_i a_i(g_i-1)$, so that $\ker \e$ is generated by $g_i-1$.
\item[(b)] We must very exactness at three stages. First we check that the sequence is exact at $\xlra{t-1}\Z G\xlra{\e}$. By the first part, we know that $\ker \e$ is generated by $t^i-1$. Now if $x\in \ker \e$, then we see that 
\begin{align*}
    x=&\sum_i a_i (t^i-1)\\
    =& \sum_i a_i (t-1)(t^{i-1}+t^{i-2}+\cdots + 1)\\
    =& (t-1)\sum_i a_i (t^{i-1}+t^{i-2}+\cdots + 1)
\end{align*}
so that $\ker \e\subseteq \im(t-1)$. Now we check that $\im(t-1)\subseteq \ker \e$. We see that $\e((t-1)x)=\e(tx-x)=\e(tx)-\e(x)=0$ and we are done. Thus the sequence is exact at this $\Z G$. 

We now check that the sequence is exact at $\xlra{N} \Z G \xlra{t-1}$.  We first note that $(t-1)(N(x))=(t-1)(\sum_i t^i x)=\sum_i t^{i+1}x-\sum_i t^ix=0$, since the two sums are equal due to $G$ being a cyclic group. For the other inclusion, if $x\in \ker (t-1)$, then we have that $tx=x$, meaning $\sum_i t^{i+1}a_i=\sum_i t^ia_i$, and by linear independence of the $t^i$'s, we have that $a_i=a_{i+1}$ for all $i$, so that $x=\sum_i a_i t^i=a_0\sum_i t^i=N(a_0)$, exactness here is proven. 

Finally, we must show that the sequence is exact at $\xlra{t-1} \Z G \xlra{N}$. We first check that $\im(t-1)\subseteq \ker N$. We see that $N((t-1)x)=\sum_i t^i (t-1)x=\sum_i t^{i+1}x-\sum_i t^ix=0$ as before. For the other inclusion, if $x\in \ker N$, then $\sum_j x t^j=0$, so that $\sum_i a_i \sum_j t^{i+j}=0$. By linear independence of the $t^{i+j} $'s, we have that $\sum_i a_i=0$, and we may write $x=\sum_i a_i t^i-\sum_i a_i=\sum_i a_i(t^i-1)$, so that $x\in \im(t-1)$ as we've shown before. This completes the proof.

\end{enumerate}

\begin{question}
\end{question}

We first note that the subgroups of $\Z G$ and $\R G$ that we are killing off are the images of the norm map $N$, so we can write $\rho_{\Z} as the quotient $\Z G/\im N$. With this in mind, we reference the previous problem to construct our resolution.
\[
0\lra \Z G/\im N \xlra{t-1} \Z G \xlra{N} \Z G/\im N \xlra{t-1} \Z G \xlra{N} \cdots
\]
and after applying the functor $\Hom_{\Z G}(-,\rho_{\Z})$, we obtain the cochain complex
\[
0\lra \Hom_{\Z G}(\Z G/\im N,\rho_{\Z}) \xlra{t-1} \Hom_{\Z G}(\Z G,\rho_{\Z}) \xlra{N} \Hom_{\Z G}(\Z G/\im N,\rho_{\Z})
\]
Noting that $\Z G$ is a rank 1 free module over itself, this cochain complex reduces to 
\[
0\lra \rho_{\Z} \xlra{t-1} \rho_{\Z} \xlra{N} \rho_{\Z}
\]
To compute $H^1(G;\rho_{\Z}), this is the same as computing $\ker (N:\rho_{\Z}\to \rho_{\Z}) / \im (t-1:\rho_{\Z}\to \rho_{\Z})$. Due to the definition of $\rho_{\Z}$, we have that $\ker N=\rho_{\Z}$, so we just need to compute the image of $t-1$. We can apply the previous problem and note that this is the same as calculating $\ker \e$. We see that $\e(x)=\e(\sum_i a_i t^i)=\sum_i a_i=0$ if and only if the sum of coefficients is divisible by $p$, and so the ideal generated by these elements is $p\Z$. Thus we have that $\im (t-1)=p\Z$, and so $H^1(G;\rho_{\Z})=\Z/p$.

The same argument applies to $\rho_{\R}$, except that since $\R$ is a field, the ideal $\langle p \rangle=\R$, so that $\im (t-1)=\R$ and thus $H^1(G;\rho_{\R})=0$.


To show that $\rho_{\Z}$ is torsion free as an abelian group, we note that as a group, it is isomorphic to the quotient $\Z^p/\Z$, where again, the $\Z$ is the subgroup generated by the image of $N$. Let $x\in \rho_{\Z}$ be a torsion element, so that $nx=0$ for some $n\in \Z$. We can write $x=\sum_i a_i t^i + \im N$, so that $nx=\sum_i na_i t^i + \im N$. Since $nx=0$, we have that $\sum_i na_i t^i\in \im N$, so that $\sum_i na_i t^i = m\sum_i t^i$ for some $m\in \Z$. By linear independence of the $t^i$'s, we have that $na_i=m$ for all $i$, so that $a_i=m/n$ for all $i$. This means that $x=m/n\sum_i t^i + \im N = m/nN(1) + \im N = 0$, so that $\rho_{\Z}$ is torsion free.

To construct the desired short exact sequence, we begin with the map $p:\rho_{\R} \to T^{p-1}_{\rho}$ given by 
\[
p(\sum_i a_i t^i+ \im N)=\sum_i e^{2\pi i a_i} t^i
\]
This map is well-defined, since if $\sum_i a_i t^i + \im N = \sum_i b_i t^i + \im N$, then $\sum_i (a_i-b_i) t^i\in \im N$, so that $\sum_i (a_i-b_i) t^i = m\sum_i t^i$ for some $m\in \Z$, and thus $e^{2\pi i a_i}=e^{2\pi i b_i}$ for all $i$. This map is also $G$-equivariant, since $p(t\cdot (\sum_i a_i t^i + \im N))=p(\sum_i a_i t^{i+1} + \im N)=\sum_i e^{2\pi i a_i} t^{i+1}=t\cdot p(\sum_i a_i t^i + \im N)$. Finally, this map is surjective, since given any $\sum_i e^{2\pi i a_i} t^i\in T^{p-1}_{\rho}$, we have that $p(\sum_i a_i t^i + \im N)=\sum_i e^{2\pi i a_i} t^i$. All that remains to be shown is that the kernel of $p$ is $\rho_{\Z}$. We first check that $\rho_{\Z}\subseteq \ker p$. If $x\in \rho_{\Z}$, then we can write $x=\sum_i a_i t^i + \im N$ where $a_i\in \Z$ for all $i$, so that $p(x)=\sum_i e^{2\pi i a_i} t^i = \sum_i t^i$, so that $x\in \ker p$. For the other inclusion, if $x\in \ker p$, then we can write $x=\sum_i a_i t^i + \im N$ where $a_i\in \R$ for all $i$, and since $p(x)=0$, we have that $\sum_i e^{2\pi i a_i} t^i = 0$, so that $e^{2\pi i a_i}=1$ for all $i$, so that $a_i\in \Z$ for all $i$, and thus $x\in \rho_{\Z}$ as desired. This completes the proof.

For the final part, note that this short exact sequence of $\Z G$-modules induces a long exact sequence in cohomology, so that we have the following exact sequence:
\[
0\lra H^0(G;\rho_{\Z}) \lra H^0(G;\rho_{\R}) \lra H^0(G;T^{p-1}_{\rho}) \lra H^1(G;\rho_{\Z}) \lra H^1(G;\rho_{\R}) 
\]
plugging in our known computations, this reduces to 
\[
H^0(G;\rho_{\R}) \lra H^0(G;T^{p-1}_{\rho}) \lra \Z/p \lra 0
\]
Now the invariants of $\rho_{\R}$ are precisely the elements in the image of $N$, which are all identified with 0, and so we see that the invariants of $T^{p-1}_{\rho}$ is isomorphic to $\Z/p$ by exactness. Since the invariants are the fixed points, this completes the proof.

We already know that $H^*(G;M)$ is well-defined, so we need to check that group cohomology is functorial with respect to maps. Let $(\a,f):(G,M)\to (G',M')$ be a morphism in $\mathcal{D}$. We define the induced map on cocahins (in the standardized cochain complex) as follows:
\[
(\a,f)^*_n\coloneqq C^n(G',M')\to C^n(G,M), \quad (\a,f)^*_n(\phi)(g_1,\ldots,g_n)=f(\phi(\a(g_1),\ldots,\a(g_n)))
\]

To see that this induces a well-defined map on cohomology, we first check that this map commutes with the differentials. We check:
\begin{align*}
&d_n((\a,f)^*_n(\phi))(g_1,\ldots,g_{n+1})\\
=& g_1\cdot (\a,f)^*_n(\phi)(g_2,\ldots,g_{n+1}) + \sum_{i=1}^n (-1)^i (\a,f)^*_n(\phi)(g_1,\ldots,g_i g_{i+1},\ldots,g_{n+1})\\
&\qquad\qquad+(-1)^{n+1} (\a,f)^*_n(\phi)(g_1,\ldots,g_n)\\
=& f(\a(g_1)\cdot \phi(\a(g_2),\ldots,\a(g_{n+1}))) + \sum_{i=1}^n (-1)^i f(\phi(\a(g_1),\ldots,\a(g_i)\a(g_{i+1}),\ldots,\a(g_{n+1})))\\
&\qquad\qquad+(-1)^{n+1}f(\phi(\a(g_1),\ldots,\a(g_n)))\\
=& f(d_n(\phi)(\a(g_1),\ldots,\a(g_{n+1})))\\
=& (\a,f)^*_{n+1}(d_n(\phi))(g_1,\ldots,g_{n+1})
\end{align*}
so that the map commutes with the differentials and we have a well-defined map on cohomology. Now we check that this map is functorial with respect to composition. Let $(\a',f'):(G',M')\to (G'',M'')$ be another morphism in $\mathcal{D}$. We check:
\begin{align*}
((\a',f')\circ (\a,f))^*_n(\phi)(g_1,\ldots,g_n) =& (\a',f')^*_n((\a,f)^*_n(\phi))(g_1,\ldots,g_n)\\
=& f'((\a,f)^*_n(\phi)(\a'(g_1),\ldots,\a'(g_n))))\\
=& f'(f(\phi(\a(\a'(g_1)),\ldots,\a(\a'(g_n)))))\\
=& (f'\circ f)(\phi((\a'\circ \a)(g_1),\ldots,(\a'\circ \a)(g_n))))\\
=& ((\a',f')\circ (\a,f))^*_n(\phi)(g_1,\ldots,g_n)
\end{align*}
so that the map is functorial with respect to composition. Finally, we check that the map is functorial with respect to identities. We check:
\begin{align*}
&(\id_G,\id_M)^*_n(\phi)(g_1,\ldots,g_n)\\
=& \id_M(\phi(\id_G(g_1),\ldots,\id_G(g_n)))\\
=& \phi(g_1,\ldots,g_n)
\end{align*}
so that the map is functorial with respect to identities as well.